\documentclass{chsmc}
\chsmcsetup{
  year = 2022,
  part = 1,
  type = A,
  show-answers = false,
}
\begin{document}

\section{填空题：本大题共 8 小题，每小题 8 分，共 64 分.}

\begin{enumerate}
  \item 集合 $A = \{n \mid n^3 < 2022 < 3^n,\; n \in \mathbf{Z}\}$
    的所有元素之和为 \blankline
  \item 设函数 $f(x) = \dfrac{x^2 + x + 16}{x}$ $(2 \leqslant x \leqslant a)$，
    其中实数 $a > 2$，若 $f(x)$ 的值域为 $[9,11]$，则 $a$ 的取值范围是 \blankline
  \item 一枚不均匀的硬币，若随机抛掷它两次均得到正面的概率是均得到反面的概率的 $9$ 倍，
    则随机抛掷它两次均得到正面、反面各一次的概率为 \blankline
  \item 若复数 $z$ 满足：$\dfrac{z - 3\ii}{z + \ii}$ 为负实数
    （$\ii$ 为虚数单位），$\dfrac{z - 3}{z + 1}$ 为纯虚数，
    则 $z$ 的值为 \blankline
  \item 若四棱锥 $P$-$ABCD$ 的棱 $AB$, $BC$ 的长均为 $\sqrt{2}$，
    其他各棱长均为 $1$，则该四棱锥的体积为 \blankline
  \item 已知函数 $y = f(x)$ 的图象既关于点 $(1,1)$ 中心对称，
    又关于直线 $x + y = 0$ 轴对称. 若 $x \in (0,1)$ 时，
    $f(x) = \log_2(x+1)$，则 $f(\log_2 10)$ 的值为 \blankline
  \item 在平面直角坐标系中，椭圆 $\Omega\colon \dfrac{x^2}{4} + y^2 = 1$，
    $P$ 为 $\Omega$ 上的动点，$A$, $B$ 为两个定点，其中 $B$ 的坐标为 $(0,3)$，
    若 $\triangle PAB$ 的面积的最小值为 $1$ 最大值为 $5$，
    则线段 $AB$ 的长为 \blankline
  \item 一个单位方格的四条边中，若有两条边染了颜色 $i$，
    另两条边分别染了异于 $i$ 色的另两种不同颜色，则称该单位方格是“$i$ 色主导”的.
    如图，一个 $1 \times 3$ 方格表的表格线共含 $10$ 条单位长线段，
    现要对这 $10$ 条线段染色，每条线段染为红、黄、蓝三色之一，
    使得红色主导、黄色主导、蓝色主导的单位方格各有一个.
    这样的染色方式数为 \blankline （答案用数值表示）
\end{enumerate}

\section{解答题：本大题共 3 个小题，满分 56 分. 解答应写出文字说明、证明过程或演算步骤.}

\begin{enumerate}[resume]
  \item (本题满分 16 分) 若 $\triangle ABC$ 的内角为 $A$, $B$, $C$ 满足
    $\sin A = \cos B = \tan C$，求 $\cos^3 A + \cos^2 A - \cos A$ 的值.
  \item (本题满分 20 分) 给定正整数 $m$ $(m \geqslant 3)$.
    设正项等差数列 $\{a_n\}$ 与正项等比数列 $\{b_n\}$ 满足：
    $\{a_n\}$ 的首项等于 $\{b_n\}$ 的公比，$\{b_n\}$ 的首项等于 $\{a_n\}$ 的公差，
    且 $a_m = b_m$. 求 $a_m$ 的最小值，并确定当 $a_m$ 取到最小值时
    $a_1$ 与 $b_1$ 的比值.
  \item (本题满分 20 分) 在平面直角坐标系中，双曲线
    $\Gamma\colon \dfrac{x^2}{3} - y^2 = 1$.
    对平面内不在 $\Gamma$ 上的任意一点 $P$. 记 $\Omega_P$ 为过点 $P$ 且与 $\Gamma$
    有两个交点的直线的全体. 对任意直线 $\ell \in \Omega_P$，
    记 $M$, $N$ 为 $\ell$ 与 $\Gamma$ 的两个交点.
    定义 $f_P(\ell) = |PM|\cdot|PN|$. 若存在一条直线
    $\ell_0 \in \Omega_P$ 满足：$\ell_0$ 与 $\Gamma$ 的两个交点位于 $y$ 轴异侧，
    且对任意直线 $\ell \in \Omega_P$，$\ell \neq \ell_0$，均有
    $f_P(\ell) > f_P(\ell_0)$，则称 $P$ 为“好点”.
    求所有好点所构成的区域的面积.
\end{enumerate}

\end{document}